\begin{textsample}{POS Dim 1 – 2013\_c4 – Score -5.00 – 2013\_c4/enade\_2013\_tec\_gest\_hosp\_e021d6e676.txt}  \label{ex:f1_pos_388}
CAPÍTULO 4 PERCEPÇÃO DA PROVA As análises feitas neste capítulo tratam das percepções dos concluintes da Área de Tecnologia em Gestão Hospitalar sobre a prova aplicada no ENADE/2013.
Estas percepções foram mensuradas por meio de nove questões que avaliaram desde o grau de dificuldade da prova até o tempo gasto para concluí -la.
As percepções sobre a prova foram relacionadas com o desempenho dos estudantes e com a Grande Região de funcionamento do curso.
O questionário de percepção da prova encontra -se ao final do Anexo IV, que traz a reprodução da prova.
O desempenho dos estudantes foi classificado em quatro quartos.
Para tanto, esse desempenho foi ordenado de forma ascendente.
O percentil 25, P25, também conhecido como primeiro quartil, é a nota de desempenho que deixa um quarto ( 25 \% ) dos valores observados abaixo e três quartos acima.
A Figura 4.1 apresenta uma ilustração deste conceito.
O quarto inferior de desempenho é composto pelas notas abaixo do primeiro quartil.
Já o percentil 75, P75, também conhecido como terceiro quartil, é o valor para o qual há três quartos ( 75 \% ) dos dados abaixo e um quarto acima dele.
O quarto superior de desempenho é composto pelas notas iguais ou acima do terceiro quartil.
O percentil 50, P50, também conhecido como mediana, é o valor que divide as notas em dois conjuntos de igual tamanho.
O segundo quarto inclui valores entre o primeiro quartil ( P25 ) e a mediana.
O terceiro quarto contém os valores entre a mediana ( P50 ) e o terceiro quartil ( P75 ).
Vale ressaltar que percentis, quartis e medianas são pontos que não obrigatoriamente pertencem ao conjunto original de dados, ao passo que os quartos são subconjuntos dos dados originais.
A seguir, serão apresentados gráficos com resultados selecionados, relativos às nove questões avaliadas por grupos de estudantes.
Os gráficos apresentam nas barras o percentual de alunos que assinalaram uma das opções ou a soma das porcentagens daqueles que assinalaram duas ( ou três ) delas.
Por exemplo, para as questões 1 e 2, os gráficos apresentam a porcentagem total de participantes que assinalaram as opções ( D ) difícil e ( E ) muito difícil.
Em cada barra foram assinalados também os extremos do intervalo de confiança de 95 \% como linhas verticais unidas por uma linha horizontal na forma da letra H maiúscula, semelhantemente aos gráficos do Capítulo 3.
As Tabelas no Anexo II apresentam os valores absolutos e a distribuição percentual12 das alternativas válidas das nove questões, segundo o mesmo recorte de desempenho dos alunos e Grande Região de funcionamento do curso. 4.1 GRAU DE DIFICULDADE DA PROVA 4.1.1 Componente de Formação Geral Ao avaliarem “ Qual o grau de dificuldade desta prova na Parte de Formação Geral? ” ( Questão 1 ), 32,0 \% do grupo de inscritos e presentes optaram pelas alternativas difícil ou muito difícil.
Entretanto, para mais da metade dos estudantes ( 60,6 \% ), o Componente de Formação Geral da prova foi considerado com grau de dificuldade médio ( Gráfico 4.1, Gráfico 4.2 e, no Anexo II, a Tabela II.1 ).
O percentual de estudantes que consideraram a prova como difícil ou muito difícil foi maior na região Sul, onde a proporção foi de 37,0 \%, enquanto a de menor incidência foi a Nordeste, com 18,3 \%.
No Gráfico 4.1 é possível observar que esta diferença é estatisticamente significativa.
Nas Grandes Regiões, a proporção de presentes à prova que consideraram o Componente de Formação Geral como sendo de grau de dificuldade médio esteve entre 57,3 \% ( regiões Sul e Centro-Oeste ) e 71,0 \% ( região Nordeste ), situação inversa dos que consideraram a prova como difícil ou muito difícil. 12 Cumpre lembrar uma das convenções para tabelas numéricas ( pág. iii ) sobre a possibilidade da soma das partes não resultar em 100 \% por questões de arredondamento.
O percentual de alunos que consideraram a prova difícil ou muito difícil foi decrescente em relação ao aumento de desempenho, do 2º ao 4º quarto ( Gráfico 4.2 ).
Nos dois quartos de menor desempenho, sem diferença estatisticamente significativa, 35,5 \% no 1º quarto e 37,2 \% no 2º consideraram a prova difícil ou muito difícil.
Nos quartos de maior desempenho a proporção de alunos que julgaram a prova difícil ou muito difícil – 29,3 \% no 3º quarto e 26,3 \% no 4º quarto – foi menor do que nos demais e não significativamente diferente entre si.
Para todos os quartos de desempenho a alternativa modal para esta pergunta foi médio, com 54,5 \%, 57,9 \%, 63,1 \% e 66,4 \% dos respondentes de cada um dos quartos, de 1 a 4 respectivamente, valores crescentes. 4.1.2 Componente de Conhecimento Específico Ao responderem à Questão 2 – “ Qual o grau de dificuldade desta prova na parte de Componente Específico? ” – 33,4 \% do grupo de estudantes classificaram -na como difícil ou muito difícil.
Além disso, o Componente de Conhecimento Específico da prova foi considerado com grau de dificuldade médio por 59,5 \% dos alunos ( Gráfico 4.3, Gráfico 4.4 e, no Anexo II, a Tabela II.2 ).
A análise das respostas dos estudantes quanto ao grau de dificuldade do Componente de Conhecimento Específico da prova, agregados por Grande Região, mostra que a diferença entre a maior e a menor proporção de alunos que a avaliaram como difícil ou muito difícil é estatisticamente significativa : 39,6 \% na região Sul e 23,8 \% na Nordeste.
O percentual de alunos que classificaram o grau de dificuldade como médio, no Componente de Conhecimento Específico, variou de 52,2 \% a 70,2 \%, para as regiões Sul e Nordeste, respectivamente, novamente situação inversa dos que consideraram a prova como difícil ou muito difícil.
Considerando -se a avaliação da dificuldade das questões do Componente de Conhecimento Específico da prova, de acordo com o desempenho dos estudantes, observa -se a não existência de diferenças estatisticamente significativas entre os quartos.
A proporção dos que classificaram a parte específica como difícil ou muito difícil variou de 29,6 \% ( 1º quarto ) a 37,4 \% ( 2º quarto ).
Já a alternativa modal para a Questão 2 foi o grau médio, com 59,3 \% do quarto inferior e 59,1 \% do superior optando por esta resposta. 4.2 EXTENSÃO DA PROVA EM RELAÇÃO AO TEMPO TOTAL Indagados quanto à extensão da prova, em relação ao tempo total oferecido para a sua resolução ( Questão 3 ), os estudantes apontaram, com maior incidência, a alternativa que considerava a extensão adequada, para todas as agregações consideradas ( Gráfico 4.5, Gráfico 4.6 e, no Anexo II, a Tabela II.3 ).
O percentual de alunos que responderam ser a extensão da prova adequada foi de 51,3 \%.
Já 41,1 \% dos inscritos presentes consideraram que o exame foi longo ou muito longo e 7,6 \% o avaliaram como curto ou muito curto.
Entre as Grandes Regiões a proporção daqueles que avaliaram a prova como longa ou muito longa em relação ao tempo total destinado à sua resolução variou de 36,7 \% na região Norte até 44,4 \% na região Nordeste.
Considerando -se o desempenho dos alunos, nota -se ainda que o percentual dos que consideraram a extensão da prova adequada variou de 49,2 \%, no 3º quarto de desempenho, a 53,5 \% no 2º quarto de desempenho.
Não se observa tendência de crescimento conforme o desempenho.
No Gráfico 4.6, pode -se constatar que para os quatro quartos de desempenho a proporção de estudantes que consideraram a prova longa ou muito longa, em relação ao tempo total destinado à sua resolução, manteve -se em torno de 41,0 \%, crescente com o desempenho, entre o 1º e o 3º quarto, e sem diferenças estatisticamente significativas. 4.3 COMPREENSÃO DOS ENUNCIADOS DAS QUESTÕES 4.3.1 Componente de Formação Geral Com relação aos enunciados das questões do Componente de Formação Geral ( Questão 4 ), as opiniões foram positivas, já que 71,6 \% dos alunos avaliados consideraram os enunciados de todas ou da maioria das questões claros e objetivos ( Gráfico 4.7, Gráfico 4.8 e, no Anexo II, a Tabela II.4 ).
Na análise regional, a percentagem de estudantes que avaliaram que todos ou a maioria dos enunciados das questões do Componente de Formação Geral estavam claros e objetivos variou de 69,7 \% nas regiões Sul e Centro-Oeste a 78,1 \% na região Nordeste, sendo que as diferenças não são estatisticamente significativas.
A análise das percepções dos estudantes sobre a clareza e objetividade dos enunciados permite afirmar que todos ou a maioria dos enunciados de questões relativas ao Componente de Formação Geral foram considerados claros e objetivos para a maior parte dos respondentes ( maior do que 69,0 \% em todas as regiões e maior do que 65,0 \% para todos os quartos de desempenho ).
Segundo o desempenho, observa -se que a proporção dos que emitiram esta opinião cresce conforme o desempenho aumenta, com diferenças estatisticamente significativas entre o 1º e o 4º quarto de desempenho, e entre o 2º e o 4º quarto de desempenho.
No quarto superior, a clareza e objetividade de todos ou da maioria dos enunciados das questões foi percebida por 78,3 \% dos alunos e no quarto de desempenho inferior tal avaliação foi emitida por 65,3 \% deles. 4.3.2 Componente de Conhecimento Específico Com relação aos enunciados das questões do Componente de Conhecimento Específico da prova, para 71,9 \% dos estudantes avaliados da Área de Tecnologia em Gestão Hospitalar a clareza e a objetividade ( Questão 5 ) estavam presentes em todas ou na maioria das questões ( Gráfico 4.9, Gráfico 4.10 e, no Anexo II, a Tabela II.5 ).
A maioria dos estudantes de todas as Grandes Regiões brasileiras considerou claros e objetivos todos ou a maioria dos enunciados das questões do Componente de Conhecimento Específico da prova, percentual sempre maior do que 69,0 \%.
As diferenças entre as regiões não são estatisticamente significativas.
A proporção de estudantes que consideraram os enunciados das questões claros e objetivos apresenta uma tendência, grosso modo, crescente em relação ao aumento de desempenho : mais elevada no quarto superior ( 78,4 \% ) se comparada ao quarto inferior de desempenho ( 66,2 \% ).
As diferenças dos resultados entre os quartos de desempenho não são estatisticamente significativas, a exceção da comparação entre o 1º e o 4º quarto. 4.4 SUFICIÊNCIA DAS INFORMAÇÕES/INSTRUÇÕES FORNECIDAS Ao avaliarem as informações/instruções fornecidas para a resolução das questões ( Questão 6 ), 77,1 \% dos respondentes da Área de Tecnologia em Gestão Hospitalar de todo o Brasil afirmaram que estas eram até excessivas ou suficientes em todas ou na maioria das questões ( Gráfico 4.11, Gráfico 4.12 e, no Anexo II, a Tabela II.6 ).
Quanto à distribuição de respondentes pelas Grandes Regiões observa -se que a proporção de estudantes que consideraram as informações/instruções fornecidas até excessivas ou suficientes em todas ou na maioria das questões foi sempre superior a 73,0 \%, chegando a 82,5 \% na região Norte.
As diferenças entre as proporções das regiões não são estatisticamente significativas.
Levando -se em conta o desempenho dos participantes, notam -se diferenças estatisticamente significativas das opiniões de estudantes entre o 1º quarto com o 3º e com o 4º quarto de desempenho, como mostra o Gráfico 4.12.
O percentual de participantes que avaliaram as informações/instruções como até excessivas ou suficientes em todas ou na maioria das questões foi mais elevado no 3º quarto ( 82,7 \% ), percentual superior à média nacional ( 77,1 \% ).
Já no quarto inferior, a suficiência das informações/instruções declarada como até excessiva em todas ou na maioria das questões foi percebida por 69,0 \% dos respondentes, com um comportamento crescente com os quartos de desempenho até o 3º quarto. 4.5 DIFICULDADE ENCONTRADA AO RESPONDER À PROVA Perguntados sobre as dificuldades com as quais se depararam ao responder à prova ( Questão 7 ), 5,4 \% dos estudantes apontaram o desconhecimento do conteúdo.
Para 54,3 \%, a forma diferente de abordagem do conteúdo foi indicada como dificuldade.
Já a falta de motivação para fazer a prova foi a dificuldade apontada por 15,0 \% dos respondentes.
Considerando -se todo o Brasil, 21,0 \% dos respondentes afirmaram que não tiveram qualquer tipo de dificuldade para responder à prova ( Tabela II.7 no Anexo II ).
Os Gráficos 4.13 e 4.14 apresentam os percentuais de estudantes que apontaram o desconhecimento do conteúdo como dificuldade percebida ao responder à prova.
Na análise por Grandes Regiões, o percentual de inscritos e presentes que apontaram o desconhecimento do conteúdo como dificuldade ao responder à prova não superou 7,0 \%.
Os percentuais variaram de 4,3 \% na região Centro-Oeste a 6,1 \% na Sul, e esta diferença não é estatisticamente significativa.
A forma diferente de abordagem do conteúdo foi a escolha modal dos estudantes, com percentuais que variaram de 48,2 \% ( região Nordeste ) a 58,9 \% ( região Sul ).
O percentual de alunos que citaram a falta de motivação como dificuldade variou de 11,5 \% ( região Sul ) a 18,6 \% ( região Centro-Oeste ).
A proporção dos que declararam não ter qualquer dificuldade para responder à prova variou de 19,8 \% na região Sul a 23,2 \% na Nordeste.
Com relação aos quartos de desempenho, o desconhecimento do conteúdo, foi a opção escolhida por 6,8 \% dos estudantes do quarto superior e 7,2 \% do quarto inferior.
A alternativa modal para os alunos, quando agregados pelos quartos de desempenho, foi que a dificuldade encontrada foi causada pela forma diferente de abordagem do conteúdo : 51,9 \% no quarto inferior e 57,5 \% do 3º quarto assim o responderam. 4.6 CONTEÚDOS DAS QUESTÕES OBJETIVAS DA PROVA Ao analisarem os conteúdos das questões objetivas da prova ( Questão 8 ), um percentual muito pequeno dos estudantes avaliados, apenas 7,5 \%, afirmou que não estudou ainda a maioria desses conteúdos ( Gráfico 4.15, Gráfico 4.16 e a Tabela II.8 no Anexo II ).
A maioria ( 66,4 \% ) afirmou ter estudado e aprendido muitos ou todos os conteúdos.
Na análise por Grande Região, a proporção de respondentes que escolheram a opção não estudou ainda a maioria desses conteúdos, foi pequena.
Observa -se que nas regiões Norte ( 8,2 \% ) e Centro-Oeste ( 10,2 \% ), apesar de pequenas, as proporções foram maiores do que a média nacional ( 7,5 \% ).
Não se observam diferenças estatisticamente significativas entre as regiões.
Em todas as regiões, a maioria dos presentes afirmou ter estudado e aprendido muitos ou todos os conteúdos, com proporções variando entre 58,5 \% na região Centro-Oeste e 75,1 \% na região Sudeste.
Considerando -se separadamente as opiniões de estudantes dos quatro quartos de desempenho, observa -se que, no quarto inferior, 12,4 \% ofereceram como resposta que não estudou ainda a maioria desses conteúdos, sendo 5,0 \% os do 3º quarto e do quarto superior com a mesma resposta, com valores decrescentes com o desempenho.
A diferença entre os alunos que optaram por este motivo de dificuldade é estatisticamente significativa entre o 1º quarto e o 3º e o 4º quarto de desempenho.
Tendo em conta o quarto superior, 77,9 \% dos alunos afirmaram ter estudado e aprendido muitos ou todos os conteúdos.
No outro extremo, no quarto inferior, 52,5 \% dos alunos responderam da mesma forma, com valores crescentes com o desempenho. 4.7 TEMPO GASTO PARA CONCLUIR A PROVA Ao responderem sobre o tempo de conclusão da prova ( Questão 9 ), mais de dois terços dos estudantes ( 73,2 \% ) afirmou ter gasto entre duas e quatro horas ( Gráfico 4.17, Gráfico 4.18 e, no Anexo II, a Tabela II.9 ).
Considerando -se as cinco Grandes Regiões brasileiras, nas regiões Sudeste ( 69,2 \% ) e Centro-Oeste ( 67,9 \% ) o percentual dos que utilizaram entre duas e quatro horas para finalizar a prova foi inferior à média nacional ( 73,2 \% ).
Em duas das cinco Grandes Regiões, Norte e Sul, o percentual de alunos que dispensaram entre duas e quatro horas para concluir a prova, superou 75,0 \%, como mostra o Gráfico 4.17.
Uma vez sendo analisadas as alternativas escolhidas pelos estudantes que se situam nos diferentes quartos de desempenho, observa -se que uma proporção bem maior de participantes no quarto superior declarou ter gasto entre duas e quatro horas para concluir a prova quando comparados com os do quarto inferior, respectivamente 79,3 \% e 63,7 \%.
As diferenças entre o 1º quarto e o 2º quarto e a entre o 1º quarto e o 4º quarto são estatisticamente significativas.
As proporções evidenciam uma tendência crescente, grosso modo, quando comparados os dois quartos extremos.

% matched lemmas: 
\end{textsample}
